\begin{exercise}
Let \(\Fld\) be the category of fields and field homomorphisms; \(\Vect\) is the category of vector spaces: objects are triples \(\Triple{K}{V}{\cdot}\) where \(K\) is a field of scalars, \(V\) is an Abelian group of vectors and \(\cdot : K \times V \to V\) is an action of scalar multiplication (which distributes both over scalar and over vector addition).
A morphism \(\Triple{K}{V}{\cdot} \to \Triple{L}{W}{\cdot}\) in \(\Vect\) is a pair \(\Pair{u}{f}\) where \(u : K \to L\) is a field homomorphism and \(f : V \to W\) a group homomorphism such that \(f(a \cdot x) = u(a) \cdot f(x)\) for all \(a \in K\) and \(x \in V\).

Check that the obvious forgetful functor \(\Vect \to \Fld\) is a fibration.
What are the fibres?
Which maps are Cartesian?
\end{exercise}

\begin{solution}
First, the fibres of \(\Vect \to \Fld\) are the categories of vector spaces over a fixed field \(K\), with morphisms being ordinary \(K\)-linear maps.

Next, we show that \(\Vect \to \Fld\) is a fibration.
Let \(u : K \to L\) be a field homomorphism, and let \(\Triple{L}{V}{\cdot}\) be an object of \(\Vect\) over \(L\).
Then we claim that \(\Pair{u}{\id{V}}\) is a Cartesian lift of \(u\).
Indeed, suppose \(v : F \to K\) is another field homomorphism and \(\Pair{w}{g} : \Triple{F}{W}{\cdot} \to \Triple{L}{V}{\cdot}\) be a morphism in \(\Vect\) over \(u \circ v\) such that the following diagram commutes.
\begin{equation*}
\begin{tikzcd}[row sep=small]
\Triple{F}{W}{\cdot} \arrow[dd, -Triangle] \arrow[dr, dashed, "?"] \arrow[drr, bend left, "{\Pair{w}{g}}"] \\
& \Triple{K}{V}{\cdot} \arrow[dd, -Triangle] \arrow[r, "{\Pair{u}{\id{V}}}"] & \Triple{L}{V}{\cdot} \arrow[dd, -Triangle] \\
F \arrow[dr, "v"] \\
& K \arrow[r, "u"] & L
\end{tikzcd}
\end{equation*}
The unique filler of the dashed ``\(?\)'' morphism in this diagram is \(? = \Pair{v}{g}\), so \(\Pair{u}{\id{V}}\) is Cartesian, and hence \(\Vect \to \Fld\) is Cartesian.

Finally, we claim that a morphism \(\Pair{u}{f} : \Triple{K}{V}{\cdot} \to \Triple{L}{W}{\cdot}\) is Cartesian if and only if \(f : V \to W\) is an abelian group isomorphism.
To show that abelian group isomorphisms induce Cartesian morphisms, mimic the proof above that showed that \(\Pair{u}{\id{V}}\) is a Cartesian lift of \(u : K \to L\) whenever \(\Triple{L}{V}{\cdot}\) is an object in \(\Vect\) above \(L\).
Conversely, suppose \(\Pair{u}{f} : \Triple{K}{V}{\cdot} \to \Triple{L}{W}{\cdot}\) is a Cartesian morphism in \(\Vect\) above a field homomorphism \(u : K \to L\).
We claim that \(f : V \to W\) is an abelian group isomorphism.
We already showed that in this case \(\Pair{u}{\id{W}} : \Triple{K}{W}{\cdot} \to \Triple{L}{W}{\cdot}\) is a Cartesian lift of \(u\), so by \expartref{1}{1}{1}{i}, there exists a unique isomorphism \(\Pair{w}{g} : \Triple{K}{V}{\cdot} \to \Triple{K}{W}{\cdot}\) in \(\Vect\) such that \(\Pair{u}{f} = \Pair{u}{\id{W}} \circ \Pair{w}{g}\).
In particular, \(g\) is an abelian group isomorphism and \(f = \id{W} \circ g = g\), so \(f\) is an abelian group isomorphism.
\end{solution}